%%%%%%%%%%%%%%%%%%%%%%%%%%%%%%%%%%%%%%%%%
% University/School Laboratory Report
% LaTeX Template
% Version 3.1 (25/3/14)
%
% This template has been downloaded from:
% http://www.LaTeXTemplates.com
%
% Original author:
% Linux and Unix Users Group at Virginia Tech Wiki 
% (https://vtluug.org/wiki/Example_LaTeX_chem_lab_report)
%
% License:
% CC BY-NC-SA 3.0 (http://creativecommons.org/licenses/by-nc-sa/3.0/)
%
%%%%%%%%%%%%%%%%%%%%%%%%%%%%%%%%%%%%%%%%%

%----------------------------------------------------------------------------------------
%	PACKAGES AND DOCUMENT CONFIGURATIONS
%----------------------------------------------------------------------------------------

\documentclass{article}
\usepackage{float}
\usepackage[version=3]{mhchem} % Package for chemical equation typesetting
\usepackage{siunitx} % Provides the \SI{}{} and \si{} command for typesetting SI units
\usepackage{graphicx} % Required for the inclusion of images
\usepackage{natbib} % Required to change bibliography style to APA
\usepackage{amsmath} % Required for some math elements 
\usepackage{fullpage} %use smaller top and bottom margins
\usepackage[utf8]{inputenc} %TODO hier ev.  latin1 setzen
%\headheight = 20pt
\headsep = 35pt %space between headline and body text
%\voffset = 0pt

\usepackage{listings}
\usepackage{courier}
\usepackage{color}   
\usepackage{titling}
\usepackage{float}

\definecolor{dkgreen}{rgb}{0,0.6,0}
\definecolor{gray}{rgb}{0.5,0.5,0.5}
\definecolor{mauve}{rgb}{0.58,0,0.82}

\setlength\parindent{0pt} % Removes all indentation from paragraphs

\renewcommand{\labelenumi}{\alph{enumi}.} % Make numbering in the enumerate environment by letter rather than number (e.g. section 6)

%\usepackage{times} % Uncomment to use the Times New Roman font



\lstset{ %
	language=R,                     % the language of the code
	basicstyle=\footnotesize,       % the size of the fonts that are used for the code
	numbers=left,                   % where to put the line-numbers
	numberstyle=\tiny\color{gray},  % the style that is used for the line-numbers
	stepnumber=1,                   % the step between two line-numbers. If it's 1, each line
	% will be numbered
	numbersep=5pt,                  % how far the line-numbers are from the code
	backgroundcolor=\color{white},  % choose the background color. You must add \usepackage{color}
	showspaces=false,               % show spaces adding particular underscores
	showstringspaces=false,         % underline spaces within strings
	showtabs=false,                 % show tabs within strings adding particular underscores
	frame=single,                   % adds a frame around the code
	rulecolor=\color{black},        % if not set, the frame-color may be changed on line-breaks within not-black text (e.g. commens (green here))
	tabsize=2,                      % sets default tabsize to 2 spaces
	captionpos=b,                   % sets the caption-position to bottom
	breaklines=true,                % sets automatic line breaking
	breakatwhitespace=false,        % sets if automatic breaks should only happen at whitespace
	title=\lstname,                 % show the filename of files included with \lstinputlisting;
	% also try caption instead of title
	keywordstyle=\color{blue},      % keyword style
	commentstyle=\color{dkgreen},   % comment style
	stringstyle=\color{mauve},      % string literal style
	escapeinside={\%*}{*)},         % if you want to add a comment within your code
	morekeywords={*,...},            % if you want to add more keywords to the set
	   literate=
	   {á}{{\'a}}1 {é}{{\'e}}1 {í}{{\'i}}1 {ó}{{\'o}}1 {ú}{{\'u}}1
	   {Á}{{\'A}}1 {É}{{\'E}}1 {Í}{{\'I}}1 {Ó}{{\'O}}1 {Ú}{{\'U}}1
	   {à}{{\`a}}1 {è}{{\`e}}1 {ì}{{\`i}}1 {ò}{{\`o}}1 {ù}{{\`u}}1
	   {À}{{\`A}}1 {È}{{\'E}}1 {Ì}{{\`I}}1 {Ò}{{\`O}}1 {Ù}{{\`U}}1
	   {ä}{{\"a}}1 {ë}{{\"e}}1 {ï}{{\"i}}1 {ö}{{\"o}}1 {ü}{{\"u}}1
	   {Ä}{{\"A}}1 {Ë}{{\"E}}1 {Ï}{{\"I}}1 {Ö}{{\"O}}1 {Ü}{{\"U}}1
	   {â}{{\^a}}1 {ê}{{\^e}}1 {î}{{\^i}}1 {ô}{{\^o}}1 {û}{{\^u}}1
	   {Â}{{\^A}}1 {Ê}{{\^E}}1 {Î}{{\^I}}1 {Ô}{{\^O}}1 {Û}{{\^U}}1
	   {œ}{{\oe}}1 {Œ}{{\OE}}1 {æ}{{\ae}}1 {Æ}{{\AE}}1 {ß}{{\ss}}1
	   {ç}{{\c c}}1 {Ç}{{\c C}}1 {ø}{{\o}}1 {å}{{\r a}}1 {Å}{{\r A}}1
	   {€}{{\EUR}}1 {£}{{\pounds}}1
} 



\setlength{\droptitle}{-10em}   % This is your set screw


%----------------------------------------------------------------------------------------
%	DOCUMENT INFORMATION
%----------------------------------------------------------------------------------------

\title{Übung 03 \\ Regression in \textbf{R} \\ INFI-IS \\ 5xHWII} 
\author{Albert Greinöcker}
\date{\today} % Date for the report

\begin{document}

\maketitle % Insert the title, author and date

\begin{center}

\includegraphics[width=10cm]{logo.png}
\end{center}
\vspace{1cm}
%----------------------------------------------------------------------------------------

\section{Bevölkerungsentwickung Tirol Gesamt}
\label{tirol_gesamt}

Die Aufgaben dieser Übung beziehen sich auf den Datensatz zur \texttt{Bevölkerungsentwicklung} aus Moodle.


Ein Modell, dass die Bevölkerungsentwicklung Tirols voraussagt, soll erstellt werden. Zu diesem Zweck müssen die Daten mit den Befehlen \texttt{aggregate} und \texttt{apply} zusammengefasst und dann mittels \texttt{lm} das Modell gebaut werden.


\section{Erste Auswertung}



\subsection{Grafische Veranschaulichung}

Eine grafische Veranschaulichung, welches die einzelnen Werte für die Gesamtbevölkerung beinhaltet, soll erstellt werden. Die Regressionsgerade soll über diese Darstellung gelegt werden.

Hier kommt das Diagramm hin:

\begin{figure}[H]
	\begin{center}
		\includegraphics[width=0.65\textwidth]{placeholder} % Include the image placeholder.png
		\caption{Bildbeschriftung 1}
	\end{center}
\end{figure}


\subsection{Prognose}

Berechne die prognostizierte Tiroler Gesamtbevölkerung im Jahr 2030, einmal direkt mit den Koeffizienten und mit dem Befehl \texttt{predict}.

\section{Bevölkerungsentwickung in der Wohngemeinde}

Erstelle dasselbe für die Gemeinde, in der man lebt. \\

\underline{Hinweis}: Es würde eigentlich reichen, das gemeinsam erstellte Skript von Moodle einfach auf diese Auswahl anzuwenden.Die Luxusvariante wäre natürlich die Verwendung von Funktionen:

\begin{lstlisting}
#Erstellen der Funktion x
x <- function (year, val)
{
	l <- lm(val~year)
	plot(year, val)
	abline(l)
}

#Aufruf der Funktion mit den Parametern y,v
y <- 2001:2015
v <- 1:15
x(y,v)
\end{lstlisting}


\section{Gegenüberstellung von Bezirken}

2 Bezirke (z.B. \texttt{IL} und \texttt{RE})  sollen bezüglich Bevölkerungswachstum gegenübergestellt werden. Bitte beides grafisch darstellen und die Steigungen beider Modelle in gesonderten Grafiken gegenüberstellen.  \\

Dass die Grafiken auch wirklich übereinander gelegt werden und alles sichtbar ist, muss man die Zeichenbereiche (mit \texttt{xlim} und \texttt{ylim}) festlegen (mit den weiteren Befehlen werden die Achsen beschriftet): \\


\section{Erkenntnisse}

Welche Erkenntnisse können nun aus den Ergebnissen über das Datenmaterial gewonnen werden? Bitte einfach zu den einzelnen Auswertungen schreiben.

\end{document}